% Generated draft: requires mathematical and source review.
\documentclass[11pt]{article}
\usepackage[margin=25mm]{geometry}
\usepackage{fontspec}
\setmainfont{Latin Modern Roman}
\usepackage{amsmath,amssymb,mathtools,unicode-math}
\setmathfont{Latin Modern Math}
\usepackage{enumitem,needspace,xurl,hyperref}
\setlength{\parindent}{0pt}
\setlength{\parskip}{5pt}
\setlength{\emergencystretch}{3em}
\newcommand{\N}{\mathbb{N}}
\newcommand{\Z}{\mathbb{Z}}
\newcommand{\Q}{\mathbb{Q}}
\newcommand{\R}{\mathbb{R}}
\newcommand{\C}{\mathbb{C}}
\newcommand{\F}{\mathbb{F}}
\newcommand{\A}{\mathbb{A}}
\newcommand{\PP}{\mathbb{P}}
\newcommand{\E}{\mathbb{E}}
\newcommand{\eps}{\varepsilon}
\newcommand{\dd}{\,\mathrm{d}}
\begin{document}
\section*{UnsolvedMath ID: 800641}
\textbf{Source problem code:} catalogue entry
\subsection{Definitions and mathematical statement}
The supplied record does not specify the ambient class of spaces or define its cover notation. The following gives a precise conditional formulation using standard cover conventions; these conventions and the intended ambient class require confirmation against the source.

Let \(X\) be a topological space. A Borel cover of \(X\) is a countable family \(\mathcal U\) of Borel subsets of \(X\) such that \(\bigcup\mathcal U=X\) and \(X\notin\mathcal U\). Write \(\mathcal B(X)\) for the family of such covers. Define the following subfamilies:
\begin{itemize}
\item \(\mathcal B_\Omega(X)\) consists of the covers \(\mathcal U\in\mathcal B(X)\) such that every finite subset of \(X\) is contained in some member of \(\mathcal U\).
\item \(\mathcal B_\Gamma(X)\) consists of the infinite covers \(\mathcal U\in\mathcal B(X)\) such that, for every \(x\in X\), the family \(\{U\in\mathcal U:x\notin U\}\) is finite.
\item \(\mathcal B_{\mathrm T}(X)\) consists of the covers \(\mathcal U\in\mathcal B(X)\) such that each point belongs to infinitely many members of \(\mathcal U\), and, for every \(x,y\in X\), at least one of the families
\[
\{U\in\mathcal U:x\in U,\ y\notin U\},\qquad
\{U\in\mathcal U:y\in U,\ x\notin U\}
\]
is finite. These are the Borel \(\tau\)-covers.
\end{itemize}
For families \(\mathcal A(X)\) and \(\mathcal C(X)\) of covers, \(\mathsf S_1(\mathcal A,\mathcal C)\) means that for every sequence \((\mathcal U_n)_{n\geq1}\) with \(\mathcal U_n\in\mathcal A(X)\), there are choices \(U_n\in\mathcal U_n\) such that \(\{U_n:n\geq1\}\in\mathcal C(X)\). The property \(\mathsf S_{\mathrm{fin}}(\mathcal A,\mathcal C)\) means that for every such sequence there are finite, possibly empty, subfamilies \(\mathcal F_n\subseteq\mathcal U_n\) such that
\[
\bigcup_{n\geq1}\mathcal F_n\in\mathcal C(X).
\]
Here selections are regarded as families of sets, without multiplicities.

For an ambient class \(\mathscr D\) of topological spaces, call a property \(P\) hereditary on \(\mathscr D\) if
\[
\forall X\in\mathscr D\;\forall Y\subseteq X,\qquad P(X)\Longrightarrow P(Y),
\]
where \(Y\) has the subspace topology and all cover families for \(Y\) are defined using Borel subsets of \(Y\).

\textbf{Problem.} With the intended ambient class \(\mathscr D\) still to be specified, is at least one of the following properties hereditary on \(\mathscr D\)?
\[
\begin{gathered}
\mathsf S_1(\mathcal B_{\mathrm T},\mathcal B_\Gamma),\qquad
\mathsf S_1(\mathcal B_{\mathrm T},\mathcal B_{\mathrm T}),\qquad
\mathsf S_1(\mathcal B_{\mathrm T},\mathcal B_\Omega),\\
\mathsf S_1(\mathcal B_{\mathrm T},\mathcal B),\qquad
\mathsf S_{\mathrm{fin}}(\mathcal B_{\mathrm T},\mathcal B_{\mathrm T}),\qquad
\mathsf S_{\mathrm{fin}}(\mathcal B_{\mathrm T},\mathcal B_\Omega).
\end{gathered}
\]
The interpretation above uses heredity with respect to every subspace. The supplied excerpt does not establish the intended ambient class or whether its surrounding context imposes any restriction on subspaces, so this remains a conditional statement schema.
\subsection{Short English statement}
The question asks whether any of six selection properties for countable Borel \(\tau\)-covers passes from a space to every subspace. The supplied excerpt leaves the ambient class and cover conventions unspecified; the displayed formulation makes those qualifications explicit.
\subsection{Sources}
\begin{itemize}
\item Local UnsolvedMath registry, permanent ID \(800641\), exact source problem code \texttt{catalogue entry}, upstream ID \(800641\), upstream revision \texttt{d371ee917ac9647fbb06cd6fc529deec6f3dda82}; supplied complete source record. The registry labels the problem \texttt{open}; this is supplied metadata, not an independently verified current mathematical status.
\item Tsaban, \emph{Selection Principles in Mathematics: A Milestone of Open Problems} (2003), Problem 14. Supplied source URL: \url{https://arxiv.org/abs/math/0312182}. The record reports access on 2026-10-16 and extraction by \texttt{arxiv\_latex}; these are supplied provenance details, not a claim of independent consultation.
\item The registry attributes the problem to Tsaban with proposed year 2003, and preserves the source context \texttt{[ideals, Problem 4]}. The bibliographic identity of \texttt{ideals} is not supplied, so the original attribution and reference cannot be resolved from this record alone.
\end{itemize}
\end{document}
